% Options for packages loaded elsewhere
% Options for packages loaded elsewhere
\PassOptionsToPackage{unicode}{hyperref}
\PassOptionsToPackage{hyphens}{url}
%
\documentclass[
  ignorenonframetext,
]{beamer}
\newif\ifbibliography
\usepackage{pgfpages}
\setbeamertemplate{caption}[numbered]
\setbeamertemplate{caption label separator}{: }
\setbeamercolor{caption name}{fg=normal text.fg}
\beamertemplatenavigationsymbolsempty
% remove section numbering
\setbeamertemplate{part page}{
  \centering
  \begin{beamercolorbox}[sep=16pt,center]{part title}
    \usebeamerfont{part title}\insertpart\par
  \end{beamercolorbox}
}
\setbeamertemplate{section page}{
  \centering
  \begin{beamercolorbox}[sep=12pt,center]{section title}
    \usebeamerfont{section title}\insertsection\par
  \end{beamercolorbox}
}
\setbeamertemplate{subsection page}{
  \centering
  \begin{beamercolorbox}[sep=8pt,center]{subsection title}
    \usebeamerfont{subsection title}\insertsubsection\par
  \end{beamercolorbox}
}
% Prevent slide breaks in the middle of a paragraph
\widowpenalties 1 10000
\raggedbottom
\AtBeginPart{
  \frame{\partpage}
}
\AtBeginSection{
  \ifbibliography
  \else
    \frame{\sectionpage}
  \fi
}
\AtBeginSubsection{
  \frame{\subsectionpage}
}
\usepackage{iftex}
\ifPDFTeX
  \usepackage[T1]{fontenc}
  \usepackage[utf8]{inputenc}
  \usepackage{textcomp} % provide euro and other symbols
\else % if luatex or xetex
  \usepackage{unicode-math} % this also loads fontspec
  \defaultfontfeatures{Scale=MatchLowercase}
  \defaultfontfeatures[\rmfamily]{Ligatures=TeX,Scale=1}
\fi
\usepackage{lmodern}

\usetheme[]{metropolis}
\ifPDFTeX\else
  % xetex/luatex font selection
\fi
% Use upquote if available, for straight quotes in verbatim environments
\IfFileExists{upquote.sty}{\usepackage{upquote}}{}
\IfFileExists{microtype.sty}{% use microtype if available
  \usepackage[]{microtype}
  \UseMicrotypeSet[protrusion]{basicmath} % disable protrusion for tt fonts
}{}
\makeatletter
\@ifundefined{KOMAClassName}{% if non-KOMA class
  \IfFileExists{parskip.sty}{%
    \usepackage{parskip}
  }{% else
    \setlength{\parindent}{0pt}
    \setlength{\parskip}{6pt plus 2pt minus 1pt}}
}{% if KOMA class
  \KOMAoptions{parskip=half}}
\makeatother


\usepackage{longtable,booktabs,array}
\usepackage{calc} % for calculating minipage widths
\usepackage{caption}
% Make caption package work with longtable
\makeatletter
\def\fnum@table{\tablename~\thetable}
\makeatother
\usepackage{graphicx}
\makeatletter
\newsavebox\pandoc@box
\newcommand*\pandocbounded[1]{% scales image to fit in text height/width
  \sbox\pandoc@box{#1}%
  \Gscale@div\@tempa{\textheight}{\dimexpr\ht\pandoc@box+\dp\pandoc@box\relax}%
  \Gscale@div\@tempb{\linewidth}{\wd\pandoc@box}%
  \ifdim\@tempb\p@<\@tempa\p@\let\@tempa\@tempb\fi% select the smaller of both
  \ifdim\@tempa\p@<\p@\scalebox{\@tempa}{\usebox\pandoc@box}%
  \else\usebox{\pandoc@box}%
  \fi%
}
% Set default figure placement to htbp
\def\fps@figure{htbp}
\makeatother





\setlength{\emergencystretch}{3em} % prevent overfull lines

\providecommand{\tightlist}{%
  \setlength{\itemsep}{0pt}\setlength{\parskip}{0pt}}



 


\setbeamerfont{title}{size=\large} \usepackage{hyperref} \setbeamertemplate{footline}[frame number] \usepackage{tikz}
\makeatletter
\@ifpackageloaded{caption}{}{\usepackage{caption}}
\AtBeginDocument{%
\ifdefined\contentsname
  \renewcommand*\contentsname{Table of contents}
\else
  \newcommand\contentsname{Table of contents}
\fi
\ifdefined\listfigurename
  \renewcommand*\listfigurename{List of Figures}
\else
  \newcommand\listfigurename{List of Figures}
\fi
\ifdefined\listtablename
  \renewcommand*\listtablename{List of Tables}
\else
  \newcommand\listtablename{List of Tables}
\fi
\ifdefined\figurename
  \renewcommand*\figurename{Figure}
\else
  \newcommand\figurename{Figure}
\fi
\ifdefined\tablename
  \renewcommand*\tablename{Table}
\else
  \newcommand\tablename{Table}
\fi
}
\@ifpackageloaded{float}{}{\usepackage{float}}
\floatstyle{ruled}
\@ifundefined{c@chapter}{\newfloat{codelisting}{h}{lop}}{\newfloat{codelisting}{h}{lop}[chapter]}
\floatname{codelisting}{Listing}
\newcommand*\listoflistings{\listof{codelisting}{List of Listings}}
\makeatother
\makeatletter
\makeatother
\makeatletter
\@ifpackageloaded{caption}{}{\usepackage{caption}}
\@ifpackageloaded{subcaption}{}{\usepackage{subcaption}}
\makeatother

\usepackage{bookmark}
\IfFileExists{xurl.sty}{\usepackage{xurl}}{} % add URL line breaks if available
\urlstyle{same}
\hypersetup{
  pdfauthor={Giorgio Arcara},
  hidelinks,
  pdfcreator={LaTeX via pandoc}}


\author{Giorgio Arcara}
\date{}

\begin{document}


\begin{frame}
% TITLE SLIDES
\title{Metodologia della Ricerca Clinica\\ \vspace{1em} \emph{07 - Simulazione e validità di performance}}
\author{Giorgio Arcara,\\ Università di Padova \\ IRCCS San Camillo, Venezia}

\titlegraphic{

\vspace*{7cm}
\includegraphics[scale=0.15]{Figures/LogoSanCamilloIRCCS_Unipd_alpha.png}
\begin{center}
\vspace{-2.2em}
\includegraphics[scale=0.15]{Figures/CC_license_3_0.png}
\end{center}
}
%\date{\today}
\maketitle
\end{frame}

\begin{frame}{Validità di performance}
\phantomsection\label{validituxe0-di-performance}
Nel contesto degli studi di forense si parla spesso di validità di
performance per valutare se una specifica valutazione è da considerare
valida

Nota che abbiamo parlato di Validità come una proprietà dei test, mentre
questa è una proprietà della valutazione

\emph{In questa specifica istanza di misurazione, sto valutando ciò che
intendo valutare?}
\end{frame}

\begin{frame}{Performance Validity Testing}
\phantomsection\label{performance-validity-testing}
\small

Esistono strumenti (pratici e concettuali) per valutare se una
performance è da considerarsi valida. Queste procedure sono note come
\textbf{Performance Validity Test (PVT)}.

Test che usano queste procedure sono anche noti come \textbf{Effort
Tests} (test dello sforzo): si indaga se il partecipante si sta
``sforzando a sufficienza''.

Ricordiamo che I test cognitivi sono test di performance massima (Vedi
slides 02. Cornice Teorica). Basta impegnarsi di meno (o ``sforzarsi di
meno'' e si invalida il risultato del test).
\end{frame}

\begin{frame}{Validità di una performance}
\phantomsection\label{validituxe0-di-una-performance}

\begin{tikzpicture}
    % Include your image
    \node[anchor=south west, inner sep=0] (image) at (0,0)
        {\includegraphics[width=0.8\textwidth]{Figures/Test_Methodology_transp.png}};

    % Set coordinate system to match the image
    \begin{scope}[x={(image.south east)}, y={(image.north west)}]
        
        % Draw a red rectangle (thin border)
        %\draw[red, thick, rounded corners] (0.25,0.3) rectangle (0.45,0.5);
        % Optional label
        %\node[red, font=\bfseries, above] at (0.35,0.5) {Region 1};
        
        %\pause
        % a red rectangle (thicker border)
        \draw[red, ultra thick] (0.605,0.65) rectangle (0.905,0.85);
        
        % text (with fine control)
        %\node[anchor=west, xshift=-8mm, yshift=20mm] at (image.east)
        %{\parbox{4cm}{\parbox{3cm}{\scriptsize \textbf{Minacce alla validità}\\ (il test non misura quello che %vorrebbe misurare)}}};
    \end{scope}
\end{tikzpicture}

\small \emph{non impegnarsi adeguatamente inficia la relazione tra costrutto da misurare, comportamento e punteggio finale}
\end{frame}

\begin{frame}{Una nota su simulazione}
\phantomsection\label{una-nota-su-simulazione}
\small

Nota che considerare la simulazione come una condizione si/no (o sto
simulando o non sto simulando) è probabilmente una semplificazione (Nies
\& Sweet, 1994)

È più ragionevole pensare che si la performance possa essere lungo un
continuum tra la performance genuina e qualle con simulazione.

Entrano inoltre in gioco alcuni aspetti (es. aspetti psicologici
inconsci) che potrebbero portare a performance deficitarie in assenza di
deficit cognitivo/danno.

Nelle slides non approfondiremo questi aspetti (molto importanti) perché
più teorici e ci concentreremo sugli aspetti statistici.
\end{frame}

\begin{frame}{Definizioni}
\phantomsection\label{definizioni}
\small

Ci sono dei concetti importanti legati che è bene esplicitare e termini
che torneranno nelle slides successive.

\textbf{Malingering (Simulazione)}: la produzione intenzionale di
problemi fisici o psicologici falsi o grossolanamente esagerati. La
motivazione del malingering è solitamente esterna (ad esempio, evitare
il servizio militare o il lavoro, ottenere un compenso finanziario,
eludere un'azione penale o procurarsi farmaci).

\textbf{Symptom Validity Assessment (SVA)}: l'insieme di tecniche
(interviste, test, raccolta evidenze) per valutare la credibilità di
sintomi (cognitivi e fisici) manifestati da un individuo.

\textbf{Performance Validity Testing (PVT)}: l'insieme di tecniche
statistiche e metodologiche per assicurarsi che il paziente stia
affrontando in maniera adeguata la situazione di valutazione (sforzo
adeguato, motivazione adeguata).
\end{frame}

\begin{frame}{Come identificare un simulatore?}
\phantomsection\label{come-identificare-un-simulatore}
Come principio generale, Si costruisicono test in maniera tale che la
performance mi dimostri che è probabile che la persona stia mentendo.
\end{frame}

\begin{frame}{Principali strategie per identificare simulatori}
\phantomsection\label{principali-strategie-per-identificare-simulatori}
\small

\textbf{1) Pattern di risposte inverosimili anche rispondendo in maniera
casuale.} Si possono usare in test a scelta forzata e vedendo se la
performance osservata è poco probabile anche rispondendo in maniera
random (e quindi suggerendo menzogna intenzionale)

\textbf{2) Test molto facili persino per pazienti deteriorati.} Valori
normativi ottenuti da popolazioni cliniche per mostrare che persino in
popolazioni clinche anche molto gravi, la performance attesa non è così
bassa.

\textbf{3) Test sviluppati con validità di criterio}: con gruppi di
persone a cui era chiesto esplicitamente di fingere un deficit, per
avere una condizione di interesse di ``simulazione''. Queste persone
possono essere sani oppure pazienti/categoria target (es. popolazioni
psichiatriche).
\end{frame}

\begin{frame}{\small 1) Pattern di risposte inverosimili, anche
rispondendo casualmente.}
\phantomsection\label{pattern-di-risposte-inverosimili-anche-rispondendo-casualmente.}
\vfill
{\center
$\text{si trovi}\ x\ \text{tale che}\ P(o < x | R) = k_s$
\vfill
}

Questa strategia può essere applicata solo nel caso di test a scelta
forzata (es. Immagine di destra o immagine di sinistra?). In questo caso
la probabilità di ottenere il pattern osservato se la scelta fosse
random (50\% per ciascuna possibilità) \(k_s\) indica un valore soglia
di probabilità molto piccolo (in genere è il solito 5\%).
\end{frame}

\begin{frame}{1) Il Test of Memory Malingering (TOMM)}
\phantomsection\label{il-test-of-memory-malingering-tomm}
Il test of Memory Malingering è il test più comune con cui è possibile
identificare questo tipo di disturbi. Si basa su un assunto molto
semplice. Nel caso di una scelta multipla a due possibilità è molto
facile modellare la probabilità di un certo numero di risposte casuali.

Consiste di due fasi: in una fase di fanno riconoscere degli item per
ricordarli (questa viene ripetuta due volte nella versione originale del
TOMM).

Nella fase test una gli item presentati nella fase di test sono
presentati in coppia con un distrattore. Per ogni coppia l'esaminato
deve indicare quale era l'item che aveva visto in precedenza.
\end{frame}

\begin{frame}{1) Il Test of Memory Malingering (TOMM)}
\phantomsection\label{il-test-of-memory-malingering-tomm-1}
\begin{center}
\includegraphics[width=0.8\linewidth,height=\textheight,keepaspectratio]{Figures/TOMM.png}
\end{center}
\center Un esempio di stimoli del TOMM
\end{frame}

\begin{frame}{1) Il Test of Memory Malingering (TOMM)}
\phantomsection\label{il-test-of-memory-malingering-tomm-2}
\small

L'assunto è semplice. Un simulatore tenderà ad indicare volontariamente
più spesso del caso l'immagine sbagliata. Questo però di fatto
suggerisce che ha riconosciuto quale è l'immagine giusta.

In caso di scelta multiple è semplice modellare la probabilità attesa di
un certo pattern di risposte (dato un certo numero di trial) usando come
riferimento la distribuzione binomiale.

es. su 50 item anche rispondendo in manira totalmente casuale ci sono
meno del 5\% di probabilità di rispondere a meno di 19 item assumendo
che c'è il 50\% di probabilità di selezionare l'item corretto anche
casualmente.

\emph{Vedi codice per esempio}
\end{frame}

\begin{frame}{1) Il Test of Memory Malingering (TOMM)}
\phantomsection\label{il-test-of-memory-malingering-tomm-3}
\small

\textbf{Vantaggi}

\begin{itemize}
\tightlist
\item
  Il TOMM è un test molto semplice e di fatto non richiede dati
  normativi. Ma solo modellare la probabilità di risposte casuali.
\item
  Si è mostra molto efficace anche nell'identificare simulatori definiti
  tramite altra procedura (buona sensibilità e specificità, vedi anche
  esempi 3) ).
\end{itemize}

\textbf{Svantaggi}

\begin{itemize}
\tightlist
\item
  Se si conosce o si intuisce come funziona il TOMM è facile aggirarlo,
  rispondendo in maniera mista o corretta e poi andare male in altri
  test.
\item
  Un cliente potrebbe essere preparato dai propri consulenti alla
  risposta a test di questo tipo.
\end{itemize}
\end{frame}

\begin{frame}{2) Test molto facili per pazienti deteriorati}
\phantomsection\label{test-molto-facili-per-pazienti-deteriorati}
\vfill
{\center
$\text{si trovi}\ x\ \text{tale che}\ P(o < x | D) = k_s$
\vfill
}

Dove \(k_s\) indica un valore di probabilità molto piccolo (es. 5\%)
persino in pazienti molto deterioriorati (\(D\)).
\end{frame}

\begin{frame}{2) Il Word Memory Test}
\phantomsection\label{il-word-memory-test}
\small

Il Word Memory Test è un test di memoria di coppie di parole (Green et
al., 1999)

È divisto in più fasi.

Nella fase 1 sono presentate 20 coppie di parole (cane/gatto,
pesce/pinna) Presentate per 6 secondi nello schermo a coppia per due
volte.

Immediatamente dopo sono presentate 40 coppie, di cui una sola parola
era presentata originariamente e l'altra è un distrattore (es.
cane/coniglio; gatto/topo).

Dopo 30 min, senza che fosse stato detto, sono presentate altre 40
coppie con altre 40 distrattori (es. cane/ratto, etc.)

\emph{Spesso in questi test (così come per ogni test di malingering) si
stressa che il test è difficile, per indurre la persona a credere che
una persona con deficit avrebbe una performance bassa (quando in realtà
non è vero).}
\end{frame}

\begin{frame}{Il Word Memory Test}
\phantomsection\label{il-word-memory-test-1}
\textbf{Vantaggi}

\begin{itemize}
\tightlist
\item
  è relativamente facile ottenere cut-off in questa maniera
\end{itemize}

\textbf{Svantaggi} - Soffre dello stesso problema di tutti I cut-off di
normalità: si fa un ragionamento inverso, ma non c'è sensibilità o
specificità nell'identificare un \emph{vero simulatore}. Per questa
ragione praticamente tutti i test, (incluso il WMT) sono anche
sviluppati utilizzando validità di criterio (vedi slides successive).
\end{frame}

\begin{frame}{3) Test sviluppati con validità di criterio}
\phantomsection\label{test-sviluppati-con-validituxe0-di-criterio}
Molti test per identificare simulatori in ambito forense seguono
principi di validità di criterio. Dove la condizione di interesse è
\emph{essere verosimilmente un simulatore}.

\(\text{Sensibilità}: P(o < y | S)\)

\(\text{Specificità}: P(o \geq  y | NS)\)

dove \(S\) indica l'essere un simulatore secondo gold standard, mentre
\(NS\) indica l'essere un non simulatore secondo gold-standard.
\end{frame}

\begin{frame}{come definire un ``simulatore''? 1/3}
\phantomsection\label{come-definire-un-simulatore-13}
\small

\begin{enumerate}
\tightlist
\item
  si creano gruppi per condizione di interesse chiedendo a persone sane
  di simulare (anche dando incentivo economico in caso di successo per
  aumentare realismo) e sono confrontati con altri sani o con pazienti
  che è inverosimile che simulino.
\end{enumerate}
\end{frame}

\begin{frame}{come definire un ``simulatore''? 2/3}
\phantomsection\label{come-definire-un-simulatore-23}
\begin{enumerate}
\setcounter{enumi}{1}
\tightlist
\item
  si identifica presenza condizione di interesse (simulatori vs non
  simulatori) sulla base di applicazione di altri strumenti di PVT
  (anche combinazione di più strumenti). In tal caso il risultato agli
  altri strumenti di PVT diventa il gold-standard.
\end{enumerate}
\end{frame}

\begin{frame}{come definire un ``simulatore''? 3/3}
\phantomsection\label{come-definire-un-simulatore-33}
\small

\begin{enumerate}
\setcounter{enumi}{2}
\tightlist
\item
  si usano gruppi in cui è estremamente verosimile che ci sia
  simulazione probable malingering cioè gruppi patologici per cui c'è
  già forte sospetto di simulazione rispetto ad altri in cui questo non
  c'è. Un classico esempio è confrontare pazienti che stanno facendo
  riabilitazione (bassa probabilità di essere simulatori) con pazienti
  che hanno chiesto valutazione per riscuotere assicurazione e che hanno
  deficit cognitivi lamentatai molto gravi rispetto a funzionamento
  quotidiano apparentemente alto (altra probabilità di essere
  simulatori). In tal caso, l'appartenza a un gruppo rispetto a questi
  criteri è il gold standard.
\end{enumerate}
\end{frame}

\begin{frame}{definizione di ``verosimile simulatore''}
\phantomsection\label{definizione-di-verosimile-simulatore}
\small

Da punto 3) di slide precedente

Slick et al.~(1999), definisce dei criteri per identificare diverse
probabilità di simulatori (franco, probabile, possibile) sulla base di
presenza di uno o più dei seguenti criteri:

\scriptsize
\emph{In determining the presence of MND (Malingered Neurocognitive Dysfunction), the case must be evaluated on the basis of four criteria: 
\begin{itemize}
\item[(A)] presence of substantial external incentive, 
\item[(B)] evidence from neuropsychological testing, 
\item[(C)] evidence from self-report, and 
\item[(D)] behaviors meeting the necessary B and C criteria are not fully accounted for by psychiatric, neurological, or developmental factors.
\end{itemize}

All diagnoses of malingering require the presence of an external incentive (Criterion A).}
\vfill
Slick, Sherman, Iverson (1999) Diagnostic Criteria for Malingered Neurocognitive
Dysfunction: Proposed Standards for Clinical Practice
and Research. \emph{The Clinical Neuropsychologist}.
\end{frame}

\begin{frame}{3) Il Reliable Digit Span}
\phantomsection\label{il-reliable-digit-span}
Il Reliable Digit Span non è altro che la somma del Digit Span Avanti e
Digit Span Indietro.

La sua utilità in ambito forense basa sul fatto che le persone possono
sovrastimare la difficoltà di questo test e avere una performance di
fatto molto (troppo bassa). Spesso però è sono riportare validità di
criterio rispetto alla possibilità di distinguere tra gruppi di
pazienti.
\end{frame}

\begin{frame}{Reliable Digit Span}
\phantomsection\label{reliable-digit-span}
\href{Figures/Sim_sens.png}{}

\begin{columns}
\column{0.5\textwidth}
\begin{figure}
\includegraphics[scale=0.35]{Figures/Sim_sens.png}
\end{figure}


\column{0.5\textwidth}
\begin{figure}
\includegraphics[scale=0.35]{Figures/Sim_sens_PPV.png}
\end{figure}
\end{columns}

\small

Mathias et al.~(2002) Detecting Malingered Neurocognitive Dysfunction
Using the Reliable Digit Span in Traumatic Brain Injury.
\emph{Assesment}.
\end{frame}

\begin{frame}{Reliable Digit Span}
\phantomsection\label{reliable-digit-span-1}
\begin{columns}
\column{0.7\textwidth}

\begin{tikzpicture}
    % Include your image
    \node[anchor=south west, inner sep=0] (image) at (0,0)
        {\includegraphics[width=0.9\textwidth]{Figures/Sim_sens_PPV.png}};

    % Set coordinate system to match the image
    \begin{scope}[x={(image.south east)}, y={(image.north west)}]
        
        % Draw a red rectangle (thin border)
        %\draw[red, thick, rounded corners] (0.25,0.3) rectangle (0.45,0.5);
        % Optional label
        %\node[red, font=\bfseries, above] at (0.35,0.5) {Region 1};
        
        
        % a red rectangle (thicker border)
        \only<2>{\draw[red, ultra thick] (0,0.38) rectangle (0.98,0.42)};
        

    \end{scope}
\end{tikzpicture}

\column{0.5\textwidth}
\scriptsize
Ricordiamo che Positive Predictive Power (o Value) indica:
Quale percentuale dei positivi sarà effettivamente un simulatore?\\
\vspace{2em}
Questo dipende da quanti “simulatori” ci aspettiamo ci siano nella popolazione.
Avere questi valori può aiutarci nell’interpretare i risultati.

\end{columns}
\pause
\scriptsize
Es. se individuo ottiene valore < 5. Potremmo dire “anche ipotizzando un base rate basso pari a .1 (cioè 10\% di simulatori nella popolazione). Una performance inferiore a 5  ha 100\% sensibilità e 92\% specificità)"
\end{frame}

\begin{frame}{3) Test sviluppati con validità di criterio.}
\phantomsection\label{test-sviluppati-con-validituxe0-di-criterio.}
\small

\textbf{Vantaggi}:

\begin{itemize}
\tightlist
\item
  Test sviluppati con validità di criterio permettono di avere valori di
  sensibilità/specificità PPV e NPV che aiutano molto
  nell'interpretazione dei risultati.
\end{itemize}

\textbf{Svantaggi}: - In questi test il limite principale è
l'adeguatezza del gruppo di ``Simulatori''. In alcuni casi si assume
siano simulatori, es da altri test. In caso in cui si chiede
explicitamente di simulare, può essere una scelta poco ecologica: nella
simulazione legata allo sviluppo del test la posta in gioco sarà
inevitabilmente più bassa (Es. Uno studente che vince 10 € in più se
finge bene non può essere paragonato ad un paziente che finge un deficit
per avere un maggiore incasso da assicurazione).
\end{frame}

\begin{frame}{Altri test usati per identificare simulatori}
\phantomsection\label{altri-test-usati-per-identificare-simulatori}
Esistono anche altri test utilizzati per identificare simultari che non
sono test di performance massima.
\end{frame}

\begin{frame}{Lo Structured Inventory of Malingered Symptoms (SIMS) 1/4}
\phantomsection\label{lo-structured-inventory-of-malingered-symptoms-sims-14}
Esistono anche altri test che non sono di performance massima che sono
legati allo studio della simulazione. Un esempio è lo Structured
Inventory of Malingered Symptoms (SIMS) che è un questionario con
risposta si/no.

È diviso in vari item raggruppati in scale che aiutano ad identificare
simulazione di disabilità intellettive, di disturbi fisici, oppure di
disturbi psichiatrici. Si ottiene anche un punteggio totale.
\end{frame}

\begin{frame}{Lo Structured Inventory of Malingered Symptoms (SIMS) 2/4}
\phantomsection\label{lo-structured-inventory-of-malingered-symptoms-sims-24}
\vfill
\small
Negli ultimi giorni faccio fatica a dormire \hfill \textbf{SI/NO}\\
\vspace{3em}
La mia memoria è talmente  peggiorata che per\\ alcuni giorni non riesco a ricordare nulla \hfill \textbf{SI/NO}

\pause

\emph{Quale è la differenza tra questi due item?}
\end{frame}

\begin{frame}{Lo Structured Inventory of Malingered Symptoms (SIMS) 3/4}
\phantomsection\label{lo-structured-inventory-of-malingered-symptoms-sims-34}
La caratteristica di alcuni items del SIMS è che definiscono alcuni
sintomi ``impossibili''.
\end{frame}

\begin{frame}{Lo Structured Inventory of Malingered Symptoms (SIMS) 4/4}
\phantomsection\label{lo-structured-inventory-of-malingered-symptoms-sims-44}
Bello studio originale SIMS ha evidenza di validità di criterio usando
una serie di gruppi di partecipanti a cui era chiesto di fingere
specifici deficit (es. Psichiatrici, cognitivi, etc.)

In particolare elevata sensibilità, specificità etc. (Smith \& Burger
1997)

\includegraphics[width=\linewidth,height=0.5\textheight,keepaspectratio]{Figures/SIMS.png}
\end{frame}

\begin{frame}{Considerazioni conclusive 1/4}
\phantomsection\label{considerazioni-conclusive-14}
\begin{center}
\includegraphics[width=0.15\linewidth,height=\textheight,keepaspectratio]{Figures/triangle.png}
\end{center}
\small L'identificazione di simulatori si basa sugli stessi principi
statistici utilizzati per identificare deficit/danno o altre condizioni
di interesse (patologia).

In questo caso il test è spesso costruito in maniera da indurre chi è
malintenzionato ad avere un certo pattern di risposte che è improbabile
e che un paziente con un reale deficit (anche grave) non commetterebbe.
\end{frame}

\begin{frame}{Considerazioni conclusive 2/4}
\phantomsection\label{considerazioni-conclusive-24}
\small

Molto spesso I test di simulazione sono associati a sensibilità e
specificità nell'identificare I simulatori e anche a PPV (Positive
Predicitive value) o NPV (Negative Predictive Value) secondo vari
scenari.

Il tema della simulazione è complesso anche perché (a fini statistici)
spesso si usa la distinzione netta simulazione/non simulazione, quando è
verosimile che siano diversi gradi di simulazione.
\end{frame}

\begin{frame}{Considerazioni conclusive 3/4}
\phantomsection\label{considerazioni-conclusive-34}
\small

Anche per I test specifici per le simulazioni le stesse considerazioni
fatte per ogni altro aspetto statistico.

Sono preferibili campioni ampi per avere delle stime più solide dei
cut-off (campioni più rappresentativi della popolazione)
\end{frame}

\begin{frame}{Considerazioni conclusive 4/4}
\phantomsection\label{considerazioni-conclusive-44}
\small

Sono presenti I limiti del confronto con un singolo campione normativo
che potrebbe non essere adeguato per i nostri scopi (es. Per età,
scolarità). Sarebbe ideale avere sempre solo persone il più simili
possibili all'individuo da valutare.

Ogni altra considerazione rimane pertinente (es. meglio norme
paese-specifiche, norme recenti, etc.)
\end{frame}




\end{document}
